\section{Introduction}\label{sec:intro}

Performance-based research funding systems evaluate institutional research capacity through statutory submission exercises, benchmarking research volume, citation density, and disciplinary concentration \cite{hicksPerformancebasedUniversityResearch2012}. In evaluative bibliometrics and university ranking frameworks, an operational assumption posits that an assessment-defined corpus aligns with the publication portfolio retrieved when querying global bibliographic databases by institutional affiliation addresses and publication year \cite{vanraanUseBibliometricAnalysis2003,vanraanFatalAttractionConceptual2005}. Discrepancies between submission inventories and index-based corpora are routinely categorized as coverage deficits, database inaccuracies, or transcription noise \cite{olenskyDataAccuracyBibliometric2015,olenskyEvaluationCitationMatching2016}.

Olensky established that bibliometric data quality is governed by three distinct structural tiers: input data accuracy at source, algorithmic entity resolution and citation matching, and corpus definition boundaries \cite{olenskyDataAccuracyBibliometric2015}. When evaluating national research corpora, corpus definition boundaries diverge along administrative lines. National research evaluation frameworks---exemplified by the UK Research Excellence Framework (REF) and Australia's Excellence in Research for Australia (ERA)---define publication eligibility by staff employment at an administrative census date rather than the institutional affiliation printed on the publication at the time of authorship \cite{mowBlackBoxResearch2010}.

In Australia's ERA 2018 assessment, 40 Higher Education Providers (HEPs) submitted 567,647 research outputs across four traditional categories (journal articles, conference papers, books, and book chapters) and non-traditional research outputs published between 2011 and 2016. Deduplication and transitive identifier clustering reconcile these submissions into 455,741 distinct canonical works. Querying the OpenAlex global bibliographic database \cite{priemOpenAlexFullyopenIndex2022} across the identical 2011--2016 window for authors affiliated with these 40 institutions retrieves 744,574 works, indicating a gross numerical divergence of 288,833 works.

This paper establishes the mathematical volume, document type composition, citation integration, and institutional attribution divergence between the assessment-defined ERA 2018 corpus and the address-defined OpenAlex corpus. We resolve three formal empirical questions:
\begin{enumerate}
    \item What is the exact volume, document type distribution, and citation density of publications carrying Australian university affiliations in OpenAlex (2011--2016) that are absent from ERA submissions?
    \item What proportion of these omitted works are directly cited by the submitted ERA corpus, and how central are they to the internal citation graph?
    \item What proportion of submitted works diverge from the submitting university in their printed affiliations, and how does this attribution divergence scale across the six-year retrospective window?
\end{enumerate}
